% Options for packages loaded elsewhere
\PassOptionsToPackage{unicode}{hyperref}
\PassOptionsToPackage{hyphens}{url}
\PassOptionsToPackage{dvipsnames,svgnames,x11names}{xcolor}
\documentclass[
  10pt,
  fontset=fandol]{ctexart}
\usepackage{xcolor}
\usepackage[margin=16mm]{geometry}
\usepackage{amsmath,amssymb}
\setcounter{secnumdepth}{-\maxdimen} % remove section numbering
\usepackage{iftex}
\ifPDFTeX
  \usepackage[T1]{fontenc}
  \usepackage[utf8]{inputenc}
  \usepackage{textcomp} % provide euro and other symbols
\else % if luatex or xetex
  \usepackage{unicode-math} % this also loads fontspec
  \defaultfontfeatures{Scale=MatchLowercase}
  \defaultfontfeatures[\rmfamily]{Ligatures=TeX,Scale=1}
\fi
\usepackage{lmodern}
\ifPDFTeX\else
  % xetex/luatex font selection
\fi
% Use upquote if available, for straight quotes in verbatim environments
\IfFileExists{upquote.sty}{\usepackage{upquote}}{}
\IfFileExists{microtype.sty}{% use microtype if available
  \usepackage[]{microtype}
  \UseMicrotypeSet[protrusion]{basicmath} % disable protrusion for tt fonts
}{}
\makeatletter
\@ifundefined{KOMAClassName}{% if non-KOMA class
  \IfFileExists{parskip.sty}{%
    \usepackage{parskip}
  }{% else
    \setlength{\parindent}{0pt}
    \setlength{\parskip}{6pt plus 2pt minus 1pt}}
}{% if KOMA class
  \KOMAoptions{parskip=half}}
\makeatother
\setlength{\emergencystretch}{3em} % prevent overfull lines
\providecommand{\tightlist}{%
  \setlength{\itemsep}{0pt}\setlength{\parskip}{0pt}}
\usepackage{amsmath,amssymb,mathtools,needspace}
\xeCJKDeclareCharClass{CJK}{"2032,"0400->"04FF,"2160->"216B,"2460->"2473,"25A0,"25C7,"25CB,"25CF}
\IfFontExistsTF{Arial Unicode MS}{\xeCJKsetup{AutoFallBack=true}\setCJKfallbackfamilyfont{\CJKrmdefault}{Arial Unicode MS}}{}
\setcounter{secnumdepth}{0}
\setlength{\emergencystretch}{2em}
\allowdisplaybreaks
\usepackage{bookmark}
\IfFileExists{xurl.sty}{\usepackage{xurl}}{} % add URL line breaks if available
\urlstyle{same}
\hypersetup{
  pdftitle={Newton 公式},
  colorlinks=true,
  linkcolor={Maroon},
  filecolor={Maroon},
  citecolor={Blue},
  urlcolor={Blue},
  pdfcreator={LaTeX via pandoc}}

\title{Newton 公式}
\author{}
\date{}

\begin{document}
\maketitle

\subsubsection{6.3.1 Newton 公式}\label{newton-ux516cux5f0f}

对于方程 \(f(x)=0\)，为要应用迭代法，必须先将它改写成 \(x=\varphi(x)\) 的形式，即需要针对所给的函数 \(f(x)\) 构造合适的迭代函数 \(\varphi(x)\)。

迭代函数 \(\varphi(x)\) 可以是多种多样的。例如，可令 \(\varphi(x)=x+f(x)\)，这时相应的迭代公式是

\href{../../source-images/pdf-163.jpeg}{原页}

\[
x_{k+1}=x_k+f(x_k).
\tag{6.3.1}
\]

一般来说，这种迭代公式不一定收敛，或者收敛的速度缓慢。

运用前述加速技巧，对于迭代过程（6.3.1），其加速公式（6.2.10）具有以下形式：

\[
\begin{cases}
\bar{x}_{k+1}=x_k+f(x_k),\\
x_{k+1}=\bar{x}_{k+1}+\dfrac{L}{1-L}(\bar{x}_{k+1}-x_k).
\end{cases}
\]

记 \(M=L-1\)，上面两个式子可以合并写成

\[
x_{k+1}=x_k-\frac{f(x_k)}{M}.
\]

这种迭代公式通常称为简化的 Newton 公式，其相应的迭代函数是

\[
\varphi(x)=x-\frac{f(x)}{M}.
\tag{6.3.2}
\]

需要注意的是，由于 \(L\) 是 \(\varphi'(x)\) 的估计值，而 \(\varphi(x)=x+f(x)\)，这里的 \(M=L-1\) 实际上是 \(f'(x)\) 的估计值。如果用 \(f'(x)\) 代替式（6.3.2）中的 \(M\)，则得到如下形式的迭代函数：

\[
\varphi(x)=x-\frac{f(x)}{f'(x)},
\]

其相应的迭代公式

\[
x_{k+1}=x_k-\frac{f(x_k)}{f'(x_k)}
\tag{6.3.3}
\]

就是著名的 Newton 公式。

\end{document}
